\section{Problem Description}
\label{sec:problem}

\begin{table}[t]
\centering
\caption{Notation used in the problem formulation.}
\label{tab:notation}
\footnotesize
\begin{tabular}{@{}l p{5.6cm}@{}}
\toprule
\textbf{Symbol} & \textbf{Description} \\
\midrule
$t = 1,\dots,N$ & Sequential task index \\
$\mathbf{s}_t$ & Static features, known at launch \\
$a_t$ & Aggregate staged input size \\
$c_t$ & Dynamic consumption signal, instrumentation-only \\
$y_t$ & True peak memory (max RSS) \\
$\hat{y}_t$ & Predicted peak memory \\
$\mathcal{H}_m(t)$ & History of prior peak-memory outcomes \\
$\mathcal{H}_c(t)$ & History of prior audited consumption traces \\
$q_t \in \{0,1\}$ & Auditing decision for task $t$ \\
$B \in (0,1]$ & Profiling budget (max audited fraction) \\
$f$ & Memory predictor \\
$\pi$ & Auditing policy \\
$L(y_t,\hat{y}_t)$ & Asymmetric loss (underestimation penalized more) \\
\bottomrule
\end{tabular}
\end{table}


A workflow executes tasks $t = 1, \dots, N$ arriving sequentially. Before launching, a predictive framework must allocate a memory estimate $\hat{y}_t$ to task $t$. Each task $t$ is associated with:

\begin{itemize}%[leftmargin=*]
\item \textbf{Static features} ($\mathbf{s}_t$): A feature vector available at prelaunch scheduling time. This vector includes workflow metadata, process type, thread count, node constraints, including the aggregate staged input size $a_t = \sum_{f \in \text{staged}(t)} \text{size}(f)$. 
\item \textbf{Dynamic Consumption} (\textbf{${c_{t}}$}):  A scalar tracking the empirical byte volume consumed by the task during runtime execution.  While highly predictive of peak memory requirements, this metric is observable exclusively through intrusive kernel-level instrumentation, rendering it causally hidden at task launch time ($c_t \notin \mathbf{s}_t$).

\item \textbf{Target Peak Memory} (\textbf{${y_t}$}): The maximum Resident Set Size (RSS) physically measured post-execution.
\end{itemize}

Because tasks run concurrently and finish asynchronously, only information
from tasks completed before this task's launch is available. We denote this
causal history by two feature vectors, constructed only from information
available before launch of task $t$:

\begin{itemize}
\item $\mathbf{h}_{m,t}$: features derived from peak-memory outcomes
$y_\tau$ of tasks completed before launch of task $t$.
\item $\mathbf{h}_{c,t}$: features derived from consumption traces $c_\tau$
of tasks completed \emph{and audited} before launch of task $t$.
\end{itemize}

We write $\mathbf{h}_t = [\mathbf{h}_{m,t}, \mathbf{h}_{c,t}]$ for the
combined history. Which prior tasks contribute to these vectors, and how
their outcomes are summarized, is a modeling choice addressed in
Section~\ref{sec:design}.


At launch, a binary decision $q_t \in \{0,1\}$ determines whether task $t$ is to be instrumented:

\begin{itemize}
\item $q_t = 1$: the task is audited, so $(c_t, y_t)$ becomes available to future $\mathcal{H}_c$ and $\mathcal{H}_m$.
\item $q_t = 0$: only $y_t$ becomes available; $c_t$ is never observed.
\end{itemize}

Auditing incurs overhead, so the fraction of audited tasks is capped by a budget $B \in (0,1]$:
$$\frac{1}{N} \sum_{t=1}^{N} q_t \le B.$$

We seek a memory predictor $f$ and an auditing policy $\pi$,


$$\begin{aligned}
\hat{y}_t = f(\mathbf{s}_t, \mathbf{h}_t), \qquad
q_t = \pi(\mathbf{s}_t, \mathbf{h}_t),
\end{aligned}$$

that jointly minimize cumulative prediction loss subject to the budget constraint:

\begin{equation}
\min_{f,\pi} \; \sum_{t=1}^{N} L(y_t, \hat{y}_t) \quad \text{s.t.} \quad \frac{1}{N}\sum_{t=1}^N \mathbb{E}[q_t] \le B
\end{equation}

$L$ is asymmetric: it penalizes underestimation ($\hat{y}_t < y_t$) more heavily than overestimation, reflecting that under-provisioning risks OOM aborts while over-provisioning only wastes reserved resources.



% \section{Problem and Running Example}

% A workflow $w$ contains named processes. A \emph{task type} is a
% workflow-process pair $(w,p)$; a \emph{task instance} $t$ invokes that type for
% one sample or interval, while version and system configuration remain context.
% At prelaunch decision time $d_t$, \toolname observes static features $S_t$,
% including staged input size
% \begin{equation}
% a_t=\sum_{f\in\mathrm{staged}(t)}\mathrm{size}(f),
% \end{equation}
% over original, intermediate, and external files. During an audited execution it
% measures
% \begin{equation}
% c_t=r_t+m_t,
% \end{equation}
% where $r_t$ sums successful explicit-read returns and $m_t$ counts attributable
% file-backed faulted pages times page size. This operational kernel signal does
% not imply that every counted byte is semantically unique.

% Let $y_t$ be peak RSS. Histories $H_m(t)$ and $H_c(t)$ summarize only tasks
% completed before $d_t$, respectively for peak memory and audited consumed data.
% Since current consumption is unavailable, CAMP predicts it before allocating:
% \begin{align}
% \hat c_t&=g(S_t,H_m(t)),\\
% \hat y_t&=f(S_t,H_m(t),H_c(t),\hat c_t),\\
% c_t,y_t&\notin\mathrm{inputs}(t).
% \label{eq:prelaunch-prediction}
% \end{align}
% At the same decision, $q_t\in\{0,1\}$ selects the eBPF or standard lane subject
% to $\sum_{t\in T_w}q_t\le B_w$. The deployment objective is to reduce
% first-attempt underallocation while controlling reservation and audit cost.

% \begin{comment}
% \begin{figure*}[t]
% \centering
% \includegraphics[width=0.98\textwidth]{source/Figures/fig_camp_live_architecture_candidate.pdf}
% \caption{CAMP's live allocation and feedback loop. Nextflow supplies prelaunch
% task features; CAMP assigns memory and selects the standard or eBPF audit lane
% in Slurm. Only completed outcomes update SQLite and later model versions.}
% \label{fig:camp-architecture}
% \end{figure*}
% \end{comment}

% \begin{table}[t]
% \caption{Temporal feature contract; prior means completed before $d_t$.}
% \label{tab:features}
% \centering
% \scriptsize
% \setlength{\tabcolsep}{2.5pt}
% \begin{tabular}{@{}p{0.11\columnwidth}p{0.55\columnwidth}p{0.25\columnwidth}@{}}
% \toprule
% Set & Fields & Availability \\
% \midrule
% $S$ & workflow/process/version/system; staged-byte components; threads/heap;
% interval/coverage & prelaunch \\
% $H_m$ & peak p50/p95/p99, spread, support/confidence, scope & prior tasks \\
% $C,H_c$ & read+mmap outcome; prior p50/p95/p99 and consumed/input ratio &
% audited outcome/prior audits \\
% $\hat c$ & current consumed-byte estimate from $S,H_m$ & prelaunch \\
% $Y$ & peak RSS, runtime, exit/OOM status & post-run only \\
% \bottomrule
% \end{tabular}
% \end{table}

% \subsection{Running Example}
% In nf-core/seqinspector, each \texttt{KRAKEN2\_KRAKEN2} invocation classifies
% one metagenomic FASTQ sample and is a task instance. Before submission,
% Nextflow exposes staged FASTQ/database sizes, process, threads, and system
% configuration. CAMP retrieves compatible completed peak and audited-consumption
% histories, predicts consumption and an upper-tail memory request, and applies
% the audit gates. Every completed task records peak RSS; the audit lane also
% records read and mmap-fault bytes. Completion updates SQLite and history
% immediately and enters the next periodic model batch, so only later,
% not-yet-scheduled tasks can use the new evidence.
